\documentclass[10pt]{article}

% ------------------------------------------------------------
% Preamble and Packages
% ------------------------------------------------------------
\usepackage[utf8]{inputenc}
\usepackage[T1]{fontenc}
\usepackage[english]{babel}
\usepackage{amsmath,amsfonts,amssymb,amsthm,mathtools}
\usepackage{mdframed}
\usepackage{graphicx}
\usepackage[center]{caption}
\usepackage{subcaption}
\usepackage{wrapfig}
\usepackage{float}
\usepackage{makeidx}
\usepackage{listings}
\usepackage{color}
\usepackage{hyperref}
\usepackage{siunitx}
\usepackage{pdfpages}
\usepackage{lscape}
\usepackage{footnote}
\usepackage{appendix}
\usepackage{sidecap}
\usepackage{tabularx}
\makesavenoteenv{tabular}
\makesavenoteenv{table}
\usepackage{array}
\usepackage{booktabs}
\usepackage{tikz}
\usetikzlibrary{shapes.geometric,calc,positioning,3d,intersections,arrows}
\usetikzlibrary{decorations.markings}
\usepackage{pgfplots}
\usepackage{circuitikz}
\usepackage[margin=0.8in]{geometry}
\usepackage{multicol}

% Listing style
\definecolor{codegreen}{rgb}{0,0.6,0}
\definecolor{codegray}{rgb}{0.5,0.5,0.5}
\definecolor{codepurple}{rgb}{0.58,0,0.82}
\definecolor{backcolour}{rgb}{0.95,0.95,0.92}

\lstdefinestyle{mystyle}{
    backgroundcolor=\color{backcolour},
    commentstyle=\color{codegreen},
    keywordstyle=\color{magenta},
    numberstyle=\tiny\color{codegray},
    stringstyle=\color{codepurple},
    basicstyle=\ttfamily\footnotesize,
    breakatwhitespace=false,
    breaklines=true,
    captionpos=b,
    keepspaces=true,
    numbers=left,
    numbersep=5pt,
    showspaces=false,
    showstringspaces=false,
    showtabs=false,
    tabsize=2
}
\lstset{style=mystyle}

\begin{document}


\title{Home Exam - FYS-3033 Deep Learning}

\maketitle

\section*{1}

\subsection*{1a)}

As introduced in Lecture 13 of FYS-3033, Generative Adversarial Networks (GANs) are designed to learn the underlying data distribution \( p_{\text{data}} \) by training two networks in opposition: 
\begin{itemize}
\item \( G \) (the generator) generates samples from a simple distribution \( p_z(z) \), such as a Gaussian distribution, and aims to produce samples that are indistinguishable from real data.
\item \( D \) (the discriminator) is a binary classifier that attempts to distinguish between real data samples and those generated by \( G \). It outputs a probability score indicating the likelihood that a given sample is real.
\end{itemize}
The original minimax objective for the discriminator is given by:
\[
V(G, D) = \int_x p_{\text{data}}(x) \log D(x) \, dx + \int_z p_z(z) \log(1 - D(G(z))) \, dz
\]
We denote \( p_g(x) \) as the distribution induced by the generator \( G(z) \), allowing us to write:
\[
V(D) = \int_x \left[ p_{\text{data}}(x) \log D(x) + p_g(x) \log(1 - D(x)) \right] dx
\]
To find the optimal discriminator \( D^*(x) \), we take the derivative of the integrand with respect to \( D(x) \) and set it to zero:
\begin{align*}
\frac{d}{dD(x)} \left[ p_{\text{data}}(x) \log D(x) + p_g(x) \log(1 - D(x)) \right] &= 0 \\
\Rightarrow \frac{p_{\text{data}}(x)}{D(x)} - \frac{p_g(x)}{1 - D(x)} &= 0
\end{align*}
Solving for \( D(x) \):
\[
D^*(x) = \frac{p_{\text{data}}(x)}{p_{\text{data}}(x) + p_g(x)}
\]
Assuming the generator has succeeded in perfectly modeling the data distribution, i.e. \( p_g = p_{\text{data}} \), we get:
\[
D^*(x) = \frac{p_{\text{data}}(x)}{p_{\text{data}}(x) + p_{\text{data}}(x)} = \frac{1}{2}
\]
This result matches what was discussed in Lecture 13, where the discriminator becomes maximally uncertain when it cannot distinguish between real and generated data.
\medskip
To find the value of the objective function at this equilibrium, we substitute \( D^*(x) = \frac{1}{2} \) back into the loss:
\begin{align*}
V(D^*) &= \int_x p_{\text{data}}(x) \log \left( \frac{1}{2} \right) dx + \int_x p_g(x) \log \left( 1 - \frac{1}{2} \right) dx \\
&= \log \left( \frac{1}{2} \right) \cdot \left[ \int_x p_{\text{data}}(x) dx + \int_x p_g(x) dx \right] \\
&= \log \left( \frac{1}{2} \right) + \log \left( \frac{1}{2} \right) = \log \left( \frac{1}{4} \right) = -\log 4
\end{align*}
Therefore, the optimal discriminator is \( \boxed{D(x) = \frac{1}{2}} \), and the value of the loss at equilibrium is \( \boxed{-\log 4} \), as discussed in the slides from Lecture 13 (GAN Training — Minimax game).
\subsection*{1b)}
The reparameterization trick, as introduced in the lecture slides on Variational Autoencoders (Lecture 12, FYS-3033), is used to express sampling from a learned distribution in a differentiable way. Instead of sampling $z$ directly from $\mathcal{N}(\mu_{z|x}, \sigma_{z|x}^2)$, we rewrite the sampling operation as:
\[
z = \mu_{z|x} + \sigma_{z|x} \cdot \epsilon, \quad \text{where } \epsilon \sim \mathcal{N}(0, 1)
\]

\noindent This allows the randomness to be pushed into $\epsilon$, enabling backpropagation through $\mu_{z|x}$ and $\sigma_{z|x}$. We now show that $z$ indeed follows a normal distribution:

\begin{align*}
\mathbb{E}[z] &= \mathbb{E}[\mu_{z|x} + \sigma_{z|x} \cdot \epsilon] = \mu_{z|x} + \sigma_{z|x} \cdot \mathbb{E}[\epsilon] = \mu_{z|x} \\
\text{Var}(z) &= \text{Var}(\sigma_{z|x} \cdot \epsilon) = \sigma_{z|x}^2 \cdot \text{Var}(\epsilon) = \sigma_{z|x}^2 \\
\Rightarrow z &\sim \mathcal{N}(\mu_{z|x}, \sigma_{z|x}^2)
\end{align*}

\noindent This matches the target distribution for the latent variable in a VAE, confirming the validity of the trick.
\subsection*{1c)}

Let \( p(x) = \text{Bern}(p_1) \) and \( q(x) = \text{Bern}(p_2) \). The KL divergence between two discrete distributions is defined as:

\[
KL(p \parallel q) = \sum_{x \in \{0,1\}} p(x) \log \left( \frac{p(x)}{q(x)} \right)
\]

For Bernoulli distributions, we have:

\begin{align*}
p(1) &= p_1, &\quad q(1) &= p_2 \\
p(0) &= 1 - p_1, &\quad q(0) &= 1 - p_2
\end{align*}

\noindent Plugging in, we get:

\begin{align*}
KL(p \parallel q) &= p(0) \log \left( \frac{p(0)}{q(0)} \right) + p(1) \log \left( \frac{p(1)}{q(1)} \right) \\
&= (1 - p_1) \log \left( \frac{1 - p_1}{1 - p_2} \right) + p_1 \log \left( \frac{p_1}{p_2} \right)
\end{align*}

We now aim to show that this is equal to:

\[
\log \left( \frac{1 - p_1}{1 - p_2} \right) + p_1 \log \left( \frac{p_1 (1 - p_2)}{p_2 (1 - p_1)} \right)
\]

Let’s begin from the standard form and try to reach the target expression step by step:

\begin{align*}
KL(p \parallel q)
&= (1 - p_1) \log \left( \frac{1 - p_1}{1 - p_2} \right) + p_1 \log \left( \frac{p_1}{p_2} \right) \\
\intertext{We factor out a full log term from the first summand:}
&= \log \left( \frac{1 - p_1}{1 - p_2} \right) \cdot (1 - p_1) + p_1 \log \left( \frac{p_1}{p_2} \right) \\
\intertext{Now write the first term as:}
&= \log \left( \frac{1 - p_1}{1 - p_2} \right) + (1 - p_1 - 1) \log \left( \frac{1 - p_1}{1 - p_2} \right) + p_1 \log \left( \frac{p_1}{p_2} \right) \\
&= \log \left( \frac{1 - p_1}{1 - p_2} \right) - p_1 \log \left( \frac{1 - p_1}{1 - p_2} \right) + p_1 \log \left( \frac{p_1}{p_2} \right) \\
\intertext{Now combine the two terms involving \( p_1 \):}
&= \log \left( \frac{1 - p_1}{1 - p_2} \right) + p_1 \left[ \log \left( \frac{p_1}{p_2} \right) - \log \left( \frac{1 - p_1}{1 - p_2} \right) \right] \\
\intertext{Use log subtraction rule \( \log a - \log b = \log \left( \frac{a}{b} \right) \):}
&= \log \left( \frac{1 - p_1}{1 - p_2} \right) + p_1 \log \left( \frac{p_1 (1 - p_2)}{p_2 (1 - p_1)} \right)
\end{align*}

\noindent This is the desired expression. \qed

\medskip

\noindent \textbf{Usage in VAEs:} \\
In VAEs, the KL divergence between the approximate posterior \( q(z|x) \) and the prior \( p(z) \) is used as a regularization term in the Evidence Lower Bound (ELBO). When using discrete latent variables such as Bernoulli distributions, this exact KL expression allows us to compute the divergence in closed form and include it directly in the loss function, helping guide the encoder toward distributions close to the prior.
This is crucial for ensuring that the latent space is well-structured and that the decoder can generate meaningful samples from the learned distribution.

\end{document}
